\documentclass[11pt]{article}

\usepackage{fontspec}
\usepackage{polyglossia}
\setdefaultlanguage{russian}
\setotherlanguage{english}

\setmainfont{Georgia}
\setsansfont{PT Sans}
\setmonofont{Menlo}[Scale=0.85]

% Georgia/PT Sans/Menlo carry Cyrillic glyphs but don't declare the OT "Cyrl"
% script tag, so polyglossia needs these explicit Cyrillic families.
\newfontfamily\cyrillicfont{Georgia}
\newfontfamily\cyrillicfontsf{PT Sans}
\newfontfamily\cyrillicfonttt{Menlo}[Scale=0.85]

\usepackage{amsmath}
\usepackage[a4paper,margin=2.2cm]{geometry}
\usepackage{tikz}
\usepackage{pgfplots}
\pgfplotsset{compat=1.17}
\usetikzlibrary{positioning,arrows.meta,backgrounds,calc,fit}
\usepackage{booktabs}
\usepackage{array}
\usepackage{xcolor}
\usepackage{microtype}

% ---- palette ----
\definecolor{accent}{HTML}{0B6FA4}
\definecolor{accent2}{HTML}{C0392B}
\definecolor{ok}{HTML}{1E8449}
\definecolor{ink}{HTML}{1A1A1A}
\definecolor{soft}{HTML}{F4F7FA}
\definecolor{softline}{HTML}{D6E1EA}
\definecolor{screen}{HTML}{2C3E50}
\definecolor{sensor}{HTML}{8E44AD}

\color{ink}

\makeatletter
\renewcommand\section{\@startsection{section}{1}{\z@}%
  {1.5ex \@plus 1ex \@minus .2ex}{0.7ex \@plus.2ex}%
  {\sffamily\Large\bfseries\color{accent}}}
\renewcommand\subsection{\@startsection{subsection}{2}{\z@}%
  {1.2ex \@plus 1ex \@minus .2ex}{0.5ex \@plus.2ex}%
  {\sffamily\large\bfseries\color{ink}}}
\renewcommand\@seccntformat[1]{\csname the#1\endcsname.\quad}
\makeatother

\newsavebox{\keybuf}
\newenvironment{keybox}
 {\begin{lrbox}{\keybuf}\begin{minipage}{\dimexpr\linewidth-2\fboxsep-2\fboxrule}}
 {\end{minipage}\end{lrbox}%
  \par\smallskip\noindent
  \setlength{\fboxsep}{9pt}\setlength{\fboxrule}{0.9pt}%
  \fcolorbox{accent}{soft}{\usebox{\keybuf}}\par\smallskip}

\renewcommand{\arraystretch}{1.25}

\begin{document}

% ====================== TITLE ======================
\begin{center}
{\sffamily\bfseries\fontsize{22}{26}\selectfont\color{accent}
XOR на свету}\\[4pt]
{\sffamily\large почему один слой не может — и как обратная связь это чинит}\\[10pt]
{\itshape\small Подробный разбор: от того, почему одно усреднённое число не способно
вычислить XOR, до пошаговой двухслойной схемы через обратную связь экран\,$\to$\,камера.}\\[2pt]
{\small Проект \textbf{«Светоч»} — оптический нейрокомпьютер на телефоне, зеркале и фронтальной камере.}
\end{center}

\vspace{0.4em}
\hrule height 0.8pt
\vspace{0.9em}

\noindent Эта заметка отвечает на два частых вопроса по эксперименту
\textbf{«Логические вентили»} (\texttt{stage4\_logic}):

\begin{enumerate}
\item Камера интегрирует по всему кадру — или есть варианты?
\item Можно ли через обратную связь экран\,$\to$\,камера проэмулировать двухслойную
сеть и получить XOR?
\end{enumerate}

\begin{keybox}
\textbf{Короткий ответ.} При считывании \textbf{одним числом} XOR невозможен в
принципе (и это честно показано в эксперименте). При \textbf{поблочном}
считывании петля экран\,$\to$\,камера превращается в полноценный второй слой —
и XOR появляется.
\end{keybox}

% ====================== 1 ======================
\section{Что такое XOR и почему он «трудный»}

Четыре логические функции от двух входов. Нас интересует последний столбец:

\begin{center}
\begin{tabular}{cc|cccc}
\toprule
$x_1$ & $x_2$ & AND & OR & NAND & \textbf{XOR}\\
\midrule
0 & 0 & 0 & 0 & 1 & \textbf{0}\\
0 & 1 & 0 & 1 & 1 & \textbf{\color{accent2}1}\\
1 & 0 & 0 & 1 & 1 & \textbf{\color{accent2}1}\\
1 & 1 & 1 & 1 & 0 & \textbf{0}\\
\bottomrule
\end{tabular}
\end{center}

Нарисуем входы как четыре точки на плоскости $(x_1,x_2)$ и раскрасим по ответу.
Слева — OR (и любой «обычный» вентиль), справа — XOR.

\begin{center}
\begin{tikzpicture}[scale=2.4,font=\small]
% ---------- LEFT: OR separable ----------
\begin{scope}
  \draw[->,softline,line width=1pt] (-0.15,0)--(1.25,0) node[right,black]{$x_1$};
  \draw[->,softline,line width=1pt] (0,-0.15)--(0,1.25) node[above,black]{$x_2$};
  % separating line x1+x2=0.5
  \draw[accent,line width=1.4pt,dashed] (0.65,-0.12)--(-0.12,0.65);
  % points
  \fill[ok] (0,0) circle (1.6pt) node[below left]{\scriptsize 0};
  \fill[accent2] (0,1) circle (1.6pt) node[above left]{\scriptsize 1};
  \fill[accent2] (1,0) circle (1.6pt) node[below right]{\scriptsize 1};
  \fill[accent2] (1,1) circle (1.6pt) node[above right]{\scriptsize 1};
  \node[align=center,font=\sffamily\bfseries] at (0.55,1.42){OR};
  \node[align=center,font=\scriptsize\itshape,accent] at (0.55,-0.42)
       {одна прямая\\ разделяет};
\end{scope}
% ---------- RIGHT: XOR not separable ----------
\begin{scope}[xshift=2.5cm]
  \draw[->,softline,line width=1pt] (-0.15,0)--(1.25,0) node[right,black]{$x_1$};
  \draw[->,softline,line width=1pt] (0,-0.15)--(0,1.25) node[above,black]{$x_2$};
  % two lines needed
  \draw[accent,line width=1.2pt,dashed] (0.35,-0.12)--(-0.12,0.35);
  \draw[accent,line width=1.2pt,dashed] (1.12,0.65)--(0.65,1.12);
  \fill[ok] (0,0) circle (1.6pt) node[below left]{\scriptsize 0};
  \fill[accent2] (0,1) circle (1.6pt) node[above left]{\scriptsize 1};
  \fill[accent2] (1,0) circle (1.6pt) node[below right]{\scriptsize 1};
  \fill[ok] (1,1) circle (1.6pt) node[above right]{\scriptsize 0};
  \node[align=center,font=\sffamily\bfseries] at (0.55,1.42){XOR};
  \node[align=center,font=\scriptsize\itshape,accent2] at (0.55,-0.42)
       {нужны \textbf{две}\\ прямые};
\end{scope}
\end{tikzpicture}
\end{center}

\begin{itemize}
\item \textbf{AND, OR, NAND}: единицы (\textcolor{accent2}{$\bullet$}) и нули
(\textcolor{ok}{$\bullet$}) разделяются \textbf{одной прямой} — это
\emph{линейно разделимые} функции.
\item \textbf{XOR}: единицы стоят по одной диагонали, нули — по другой.
\textbf{Никакая одна прямая} их не разделит (Минский–Пейперт, 1969).
\end{itemize}

\noindent\textbf{Запомните:} «одна прямая» $=$ «один линейный слой» $=$ «одно
решающее правило». XOR требует \textbf{двух} прямых (двух скрытых нейронов) и ещё
одного правила поверх них.

% ====================== 2 ======================
\section{Как «Светоч» считает на свету}

Сенсор камеры за время экспозиции \textbf{складывает падающие фотоны}. Если на
экране показать яркости, то яркость на сенсоре — это физическая сумма (а с
весами — скалярное произведение). Сложение и скалярное произведение — базовая
операция любой нейросети, и здесь она происходит «бесплатно», за один оптический
такт. Ключевой вопрос — \textbf{как мы это сложение считываем}.

% ====================== 3 ======================
\section{Почему текущий XOR-эксперимент честно проваливается}

В \texttt{stage4\_logic.js} выбрана \textbf{самая простая} схема считывания:
оба входа кодируются \emph{суммарной яркостью} $(x_1+x_2)/2$ на весь экран, а
функция \texttt{regionMean} усредняет центральную область кадра (25–75\,\% по
обеим осям) в \textbf{один скаляр}.

\begin{center}
\begin{minipage}{0.46\textwidth}
\begin{tikzpicture}
\begin{axis}[
  width=\textwidth,height=5.2cm,
  ybar,bar width=14pt,
  ymin=0,ymax=1.15,
  ylabel={яркость на сенсоре},
  ylabel style={font=\footnotesize},
  symbolic x coords={(0{,}0),(0{,}1),(1{,}0),(1{,}1)},
  xtick=data,
  x tick label style={font=\footnotesize},
  ytick={0,0.5,1},
  nodes near coords,
  nodes near coords style={font=\scriptsize},
  every axis plot/.append style={draw=accent,fill=accent!25},
  axis lines=left,
]
\addplot coordinates {((0{,}0),0) ((1{,}1),1)};
\addplot[draw=accent2,fill=accent2!30] coordinates {((0{,}1),0.5) ((1{,}0),0.5)};
\end{axis}
\end{tikzpicture}
\end{minipage}\hfill
\begin{minipage}{0.5\textwidth}
\small
\textbf{Коллапс яркости.} Вход $(x_1,x_2)$ превращается в одно число $(x_1+x_2)/2$:
\smallskip

\begin{tabular}{c|c}
$(x_1,x_2)$ & яркость \\
\midrule
$(0,0)$ & $0.0$ \\
$(0,1)$ & \textcolor{accent2}{$\mathbf{0.5}$} \\
$(1,0)$ & \textcolor{accent2}{$\mathbf{0.5}$} \\
$(1,1)$ & $1.0$ \\
\end{tabular}
\smallskip

Случаи \textcolor{accent2}{$(0,1)$ и $(1,0)$ дают одинаковый отсчёт $0.5$}.
Камера видит «сколько всего света», а не «откуда» — и различить их \emph{не может}.
\end{minipage}
\end{center}

\noindent А ведь XOR требует, чтобы «средние» яркости ($0.5$) дали \textbf{1}, а
«крайние» ($0.0$ и $1.0$) — \textbf{0}. Это два порога (снизу и сверху), один
порог так не умеет. Поэтому эксперимент \textbf{специально} объявляет успехом
\emph{низкую} точность XOR — провал здесь запланирован и доказывает теорему
Минского–Пейперта на реальном железе, а не в симуляции. AND/OR/NAND при этом
проходят на 100\,\%.

\begin{keybox}
\textbf{Вывод по вопросу №1.} Да, в этом эксперименте камера интегрирует по всему
(центральному) кадру в одно число. Но это \textbf{выбор схемы считывания}, а не
предел железа. И при таком считывании обратная связь действительно почти
бесполезна — комментатор прав.
\end{keybox}

% ====================== 4 ======================
\section{Вариант есть: поблочное считывание}

Камера \textbf{не обязана} давать одно число. В \texttt{index.html} есть
инструменты пространственного показа и считывания:

\begin{itemize}
\item \texttt{showBlockPattern(values, gridSize)} — показать на экране
\textbf{сетку} разных яркостей (а не одну заливку);
\item \texttt{blockMeans(frame, gridSize)} — считать кадр как \textbf{сетку
регионов} и вернуть массив средних по каждой клетке.
\end{itemize}

\begin{center}
\begin{tikzpicture}[font=\footnotesize]
  % screen grid
  \foreach \x in {0,1,2,3}{\foreach \y in {0,1,2,3}{
    \pgfmathsetmacro{\b}{20+18*mod(\x+2*\y,5)}
    \fill[screen!\b] (\x*0.42,\y*0.42) rectangle ++(0.4,0.4);
    \draw[white,line width=0.6pt] (\x*0.42,\y*0.42) rectangle ++(0.4,0.4);
  }}
  \node[below=2pt,font=\sffamily\bfseries\footnotesize] at (0.84,0) {ЭКРАН: сетка яркостей};
  % arrow
  \draw[-{Stealth[length=4mm]},accent,line width=1.4pt] (2.1,0.84)--(3.2,0.84)
       node[midway,above,font=\scriptsize]{свет};
  % sensor grid with numbers
  \begin{scope}[xshift=3.4cm]
  \foreach \x in {0,1,2,3}{\foreach \y in {0,1,2,3}{
    \draw[sensor,line width=0.7pt,fill=sensor!10] (\x*0.42,\y*0.42) rectangle ++(0.4,0.4);
  }}
  \node[below=2pt,font=\sffamily\bfseries\footnotesize] at (0.84,0) {СЕНСОР: массив скаляров};
  \end{scope}
\end{tikzpicture}
\end{center}

\noindent Так уже работают «взрослые» эксперименты — скалярное произведение
(\texttt{stage2}), MatVec (\texttt{stage3}) и слой трансформера
(\texttt{stage5}): разные участки экрана и сенсора \textbf{параллельно} считают
разные скалярные произведения — десятки независимых «нейронов» за один кадр. Как
только мы умеем считывать \textbf{несколько чисел}, у нас появляется несколько
скрытых нейронов — ровно то, чего не хватало для XOR.

% ====================== 5 ======================
\section{XOR как двухслойная сеть через обратную связь}

XOR раскладывается в классику:
\[
\text{XOR}(x_1,x_2)=\text{AND}\big(\,\text{OR}(x_1,x_2),\ \text{NAND}(x_1,x_2)\,\big)
\]

\begin{center}
\begin{tabular}{cc|ccc}
\toprule
$x_1$ & $x_2$ & OR & NAND & AND(OR,NAND) $=$ XOR\\
\midrule
0 & 0 & 0 & 1 & 0\\
0 & 1 & 1 & 1 & \textbf{\color{accent2}1}\\
1 & 0 & 1 & 1 & \textbf{\color{accent2}1}\\
1 & 1 & 1 & 0 & 0\\
\bottomrule
\end{tabular}
\end{center}

\noindent OR и NAND \textbf{по отдельности линейно разделимы} (их «Светоч» уже
берёт на 100\,\%). Значит, нужен \textbf{слой 1} (посчитать OR и NAND),
\textbf{нелинейность} между слоями и \textbf{слой 2} (посчитать AND). Оптически —
два такта с обратной связью экран\,$\to$\,камера.

\subsection*{Архитектура: два оптических такта}

\begin{center}
\begin{tikzpicture}[font=\footnotesize,
  box/.style={rounded corners=3pt,draw,line width=1pt,minimum height=1.1cm,
              minimum width=2.2cm,align=center},
  scr/.style={box,fill=screen!12,draw=screen},
  sen/.style={box,fill=sensor!12,draw=sensor},
  ->/.default,
  ar/.style={-{Stealth[length=3mm]},line width=1.2pt},
]
% layer 1
\node[scr] (s1) {ЭКРАН · такт 1\\\scriptsize зона OR $\;|\;$ зона NAND};
\node[sen,right=1.5cm of s1] (d1) {СЕНСОР\\\scriptsize $h_1{=}$OR\ \ $h_2{=}$NAND};
\draw[ar,accent] (s1)--(d1) node[midway,above,font=\scriptsize]{свет};
% nonlinearity
\node[below=0.7cm of d1,box,fill=ok!10,draw=ok,minimum width=2.6cm] (nl)
     {нелинейность\\\scriptsize отклик сенсора $+$ порог};
\draw[ar,sensor] (d1)--(nl);
% layer 2
\node[scr,below=0.7cm of s1,minimum width=2.6cm] (s2)
     {ЭКРАН · такт 2\\\scriptsize зона AND};
\draw[ar,accent2] (nl) -- node[above,font=\scriptsize]{$h_1,h_2$ обратно на экран} (s2);
% output
\node[sen,left=1.5cm of s2,minimum width=2.0cm] (out)
     {ВЫХОД\\\scriptsize $y=$ XOR};
\draw[ar,accent] (s2)--(out) node[midway,above,font=\scriptsize]{свет};
% loop label
\node[font=\scriptsize\itshape,accent2,align=center] at ($(nl)!0.5!(s2)+(0,0.0)$){};
\end{tikzpicture}
\end{center}

\textbf{Слой 1 — два скрытых нейрона в двух зонах (один такт).} Через
\texttt{showBlockPattern} показываем входы так, что левая половина экрана
настроена на OR, правая — на NAND (разные веса/смещения). Камера через
\texttt{blockMeans} считывает \emph{два} числа $h_1,h_2$.

\textbf{Нелинейность — даётся самим железом.} Между слоями обязательно нужна
нелинейность (иначе два линейных слоя схлопываются в один). Её \emph{не нужно
добавлять отдельно} — её дают отклик сенсора (гамма $+$ насыщение, кривая
сигмоидной формы) и порог $\texttt{opt} > \texttt{thresh}$, который в
\texttt{stage4} и так применяется.

\begin{center}
\begin{tikzpicture}
\begin{axis}[
  width=8.5cm,height=4.4cm,
  xlabel={сумма света (вход слоя)}, ylabel={отсчёт сенсора},
  xlabel style={font=\footnotesize},ylabel style={font=\footnotesize},
  xmin=0,xmax=1,ymin=0,ymax=1,
  xtick={0,0.5,1},ytick={0,0.5,1},
  tick label style={font=\scriptsize},
  axis lines=left,
  domain=0:1,samples=80,
]
\addplot[accent,line width=1.6pt]{1/(1+exp(-11*(x-0.5)))};
\draw[dashed,softline] (axis cs:0.5,0)--(axis cs:0.5,0.5)--(axis cs:0,0.5);
\node[font=\scriptsize,accent,align=center] at (axis cs:0.74,0.30){сигмоидоподобный\\[-2pt]отклик};
\end{axis}
\end{tikzpicture}
\end{center}

\textbf{Слой 2 — обратная связь (второй такт).} Считанные $h_1$ (OR) и $h_2$
(NAND) \textbf{показываем обратно на экране} — это и есть петля
камера\,$\to$\,экран — и настраиваем зону на AND. Камера считывает финальный
скаляр; порог по нему даёт $0$ для $(0,0)$ и $(1,1)$, $1$ для $(0,1)$ и $(1,0)$.
\textbf{Это XOR, вычисленный на свету двумя оптическими проходами:} двухслойный
перцептрон, развёрнутый во времени.

% ====================== 6 ======================
\section{Две формы обратной связи в проекте}

\begin{itemize}
\item \textbf{Пространственно-временная петля} (описана выше): посчитали в
регионах $\to$ считали $\to$ подали назад $\to$ ещё слой. Даёт \emph{глубину}.
\item \textbf{Послесвечение OLED как память} (\texttt{stage37\_lstm}): пиксели
OLED и тракт «экран$\to$камера» физически удерживают остаточный образ после
погасания — ячейка памяти с «забыванием» (forget gate) и постоянной времени
$\tau$. Настоящая рекуррентность, где состояние хранится в затухающем свете.
\item \textbf{Camera-in-the-loop обучение} (статья VI): замыкание петли
экран$\to$камера позволяет не только считать, но и \emph{обучать} физический
оптический канал, подстраивая паттерны по тому, что увидел сенсор.
\end{itemize}

\noindent «Канал обратной связи», про который спрашивал комментатор, — это не
экзотика, а несущая идея: статичный канал считает один слой, а \textbf{замкнутый}
канал даёт глубину (пространственную) и память (временную).

% ====================== 7 ======================
\section{Итоги одним абзацем}

\begin{keybox}
\begin{itemize}
\item В демо-XOR камера усредняет кадр в \textbf{одно число} $\Rightarrow$ $(0,1)$
и $(1,0)$ неразличимы $\Rightarrow$ XOR честно проваливается (Минский–Пейперт на
железе).
\item Это ограничение \textbf{схемы считывания}, а не оптики: поблочное
считывание даёт \textbf{много} скаляров — несколько скрытых нейронов за такт.
\item $\text{XOR}=\text{AND}(\text{OR},\text{NAND})$. OR и NAND берутся слоем 1 в
двух зонах; нелинейность даёт сенсор $+$ порог; AND берётся слоем 2.
\item \textbf{Обратная связь экран\,$\to$\,камера} $=$ подать результат слоя 1
обратно на экран для слоя 2. При считывании одним числом — бесполезна; при
поблочном — это полноценный второй слой, и XOR появляется.
\end{itemize}
\end{keybox}

\vspace{0.6em}
\noindent\footnotesize\textit{Примечание.} Двухслойная версия — это проектное
предложение (следующий эксперимент): в коде \texttt{stage4\_logic.js} сейчас один
проход с \texttt{regionMean}. Реализация $=$ \texttt{showBlockPattern} на двух
зонах $+$ \texttt{blockMeans} для считывания $h_1,h_2$ $+$ второй проход с подачей
$h_1,h_2$ обратно.

\end{document}
